\begin{proposition}[7.2.4]
\label{IV.2.s7.2.4-fr}
Sea $A$ un anillo local noetheriano y pongamos
$X=\operatorname{Spec}(A)$. Las condiciones siguientes son
equivalentes:
\begin{enumerate}
\item[\textup{a)}] Para todo anillo cociente íntegro $B$ de $A$, el anillo
$B^{(1)}$ es una $B$-álgebra finita.
\item[\textup{b)}] Para todo $\mathscr O_X$-módulo coherente $\mathscr F$ y
todo subconjunto $Z$ estable por especialización, tal que para todo
$x\in\operatorname{Ass}(\mathscr F)\cap(X\setmin Z)$ se tenga
$\operatorname{codim}(\overline{\{x\}}\cap Z,\overline{\{x\}})\geqslant2$, el $\mathscr O_X$-módulo
$\mathscr H^0_{X/Z}(\mathscr F)$ es coherente.
\item[\textup{c)}] Para todo subconjunto cerrado $T$ de $X$ y todo
$\mathscr O_U$-módulo coherente $\mathscr G$ (donde $U=X\setmin T$)
tal que para todo $x\in\operatorname{Ass}(\mathscr G)$ se tenga
$\operatorname{codim}(\overline{\{x\}}\cap T,\overline{\{x\}})\geqslant2$, el módulo $i_*(\mathscr G)$, donde
$i:U\to X$ es la inclusión canónica, es un
$\mathscr O_X$-módulo coherente.
\item[\textup{d)}] Para todo anillo cociente íntegro $B$ de $A$ y todo ideal
$\mathfrak J$ de altura $\geqslant2$ en $B$, el anillo
$\displaystyle\bigcap_{\mathfrak p\not\supset\mathfrak J}B_{\mathfrak p}$ es una $B$-álgebra finita.
\item[\textup{e)}] Para todo anillo cociente íntegro $B$ de $A$ y todo anillo
local $C=B_{\mathfrak q}$ en un ideal primo $\mathfrak q$ de $B$
tal que $\dim(C)\geqslant2$, el anillo $C^{(\omega)}$ es una
$C$-álgebra finita \textup{(o, lo que es lo mismo,} si
$Y=\operatorname{Spec}(C)$, si $U$ es el complemento en $Y$ del
punto cerrado de $C$ y si $j:U\to Y$ es la inclusión canónica,
$j_*(\mathscr O_U)$ es un $\mathscr O_Y$-módulo coherente).
\item[\textup{f)}] Para todo anillo cociente íntegro $B$ de $A$, todo
anillo local $C=B_{\mathfrak q}$ en un ideal primo $\mathfrak q$
de $B$ tal que $\dim(C)\geqslant2$, y todo ideal
$\mathfrak r\in\operatorname{Ass}(\widehat C)$, se tiene
$\dim(\widehat C/\mathfrak r)\geqslant2$.
\end{enumerate}
\end{proposition}
